\documentclass[
]{article}
\usepackage{amsmath,amssymb}
\usepackage{lmodern}
\usepackage{iftex}
\author{}
\date{}

\begin{document}

\section{Question 2}

\subsection{Part 3}

In order to determine whether or not the mean $\mu$ of the first-difference
series of \textit{`PASSENGERS CARRIED'} has changed in the pre-Covid and post-Covid era, we can utilize a statistical hypothesis test such as the \textbf{t-test for the difference in means} between two independent samples. 

\subsubsection{Hypotheses}

The hypotheses are as follows:

\begin{itemize}
\item
  \textbf{Null Hypothesis:} There is no difference in the means of the
  first differenced series of passenger numbers between the two periods.
    \[
    H_0: \mu_{\text{pre}} = \mu_{\text{post}}
    \]
\item
  \textbf{Alternative Hypothesis:} There is a difference in the means of
  the first differenced series of passenger numbers between the two
  periods.
    \[
    H_a: \mu_{\text{pre}} \neq \mu_{\text{post}}
    \]
\end{itemize}

\subsubsection{Test Statistic}

The test statistic for the independent samples t-test is calculated as follows:
\[
t = \frac{\bar{X}_{\text{pre}} - \bar{X}_{\text{post}}}{\sqrt{\frac{s_{\text{pre}}^2}{n_{\text{pre}}} + \frac{s_{\text{post}}^2}{n_{\text{post}}}}}
\]
Where:
\begin{itemize}
    \item \( \bar{X}_{\text{pre}} \) and \( \bar{X}_{\text{post}} \) are the sample means for pre-COVID and post-COVID, respectively.
    \item \( s_{\text{pre}}^2 \) and \( s_{\text{post}}^2 \) are the sample variances for the two groups.
    \item \( n_{\text{pre}} \) and \( n_{\text{post}} \) are the sample sizes for the two groups.
\end{itemize}

\subsubsection{Interpretation of Results}

The results of the t-test are interpreted as follows:

\begin{itemize}
\def\labelenumi{\roman{enumi}.}
\item
  The t-statistic is compared to a critical value from the
  t-distribution table based on our significance level (commonly $\alpha = 0.05$) and the degrees of freedom.
\item
 If $t$ is greater than the critical value, we reject the null hypothesis, suggesting a   significant difference in the mean number of passengers carried before and after the COVID period.
\item
  Rejecting the null hypothesis indicates a significant difference in
  the mean number of passengers carried between the two periods.
  Conversely, failing to reject the null hypothesis suggests no
  significant difference in passenger numbers before and after the
  pandemic exists.
\end{itemize}

\subsubsection{Assumptions}

The assumptions of the independent samples t-test are as follows:

\begin{itemize}
\item
  \textbf{Independence:} The observations in each group must be
  independent of one another.
\item
  \textbf{Normality:} The sample means must be approximately normally
  distributed. Since our first-difference series is weakly stationary,
  this assumption is reasonable under the Central Limit Theorem if the
  sample sizes are large enough.
\item
  \textbf{Homogeneity of Variances:} The variances of the two groups
  should be equal (or approximately equal). This can be tested using
  Levene's test for equality of variances.
\end{itemize}

\end{document}
